% !TeX spellcheck = en_US
\documentclass[french]{yLectureNote}

\title{Électromagnétisme 3}
\subtitle{Physique}
\author{Paulhenry Saux}
\date{\today}
\yLanguage{Français}
%jesse.groenen@cemes.fr
\professor{M.Brut}
\usepackage{graphicx}%----pour mettre des images
\usepackage[utf8]{inputenc}%---encodage
\usepackage{geometry}%---pour modifier les tailles et mettre a4paper
%\usepackage{awesomebox}%---pour les boites d'exercices, de pbq et de croquis ---d\'esactiv\'e pour les TP de PC
\usepackage{tikz}%---pour deiffner + d\'ependance de chemfig
% \usepackage{tabularx}%---pour dimensionner automatiquement les tableaux avec variable X
\usepackage{awesomebox}%---Pour les boites info, danger et autres
\usepackage{menukeys}%---Pour deiffner les touches de Calculatrice
\usepackage{fancyhdr}%---pour les en-t\^ete personnalis\'ees
\usepackage{blindtext}%---pour les liens
\usepackage{hyperref}%---pour les liens (\`a mettre en dernier)
\usepackage{caption}%---pour la francisation de la l\'egende table vers Tableau
\usepackage{pifont}
\usepackage{array}%---pour les tableaux
\usepackage{lipsum}
\usepackage{yFlatTable}
\usepackage{multicol}
\usepackage{tabularx}
\usepackage{booktabs}
\newcommand{\Lim}[1]{\lim\limits_{\substack{#1}}\:}
\renewcommand{\vec}{\overrightarrow}
\newcommand{\N}[0]{\mathbb{N}}
\newcommand{\dd}{\mathrm{d}}
\newcommand{\norm}[1]{||\vec{#1}||}
\newcommand{\fo}{\psi(\vec{r},t)}
\newcommand{\foe}{\psi(\vec{r},t)\*}
\newcommand{\HH}{\hat{H}}
\newcommand{\hb}{\hbar}
\newcommand{\lap}{\nabla^2}
\newcommand{\lapcc}{\frac{\partial^2 }{\partial x^2}+\frac{\partial^2 }{\partial y^2}+\frac{\partial^2 }{\partial z^2}}
\newcommand{\E}{\vec{E}(M)}
\newcommand{\B}{\vec{B}(M)}
\newcommand{\V}{V(M)}
\newcommand{\er}{\vec{e_{\rho}}}
\newcommand{\ep}{\vec{e_{\varphi}}}
\newcommand{\pip}{\pi^+}
\newcommand{\pim}{\pi^-}
\newcommand{\mv}{\mathcal{V}}
\DeclareMathOperator{\Div}{Div}
\newcommand{\rot}{\vec{\mathrm{rot}}}
\newcommand{\grad}{\vec{\mathrm{grad}}}
%\setcounter{chapter}{3}
\begin{document}
%\chapter{Puissance et énergie du champ Em}
\setcounter{chapter}{3}
\chapter{Polarisation milieux diélectriques : aspects microscopiques}

Un conducteur parfait est équivalent à un milieu infiniment polarisable : $\chi_e \to \infty$. L'écrantage est total : $\vec{E_{int}} = \vec{0}$.
 C'est équivalent en polarisation alors qu'il n'y a pas de dipôle élémentaire. C'est un réservoir illimité d'électrons libres.
\section{Polarisabilité}
\subsection{Définition}
On se place en échelle microscopique
\begin{definition}
    $$\vec{p_j} = \epsilon_0 \alpha_j \vec{e_{loc}(M_j)}$$ avec $\vec{p_j}$ le moment dipolaire électrique induit sur l'élément j , $\vec{e_{loc}}$ le champ électrique micorcopique perçu au site $j$ par le constituant élémentaire
    et $\alpha_j$ la Polarisabilité de l'élément j.
\end{definition}
\warningInfo{Champ local}{
    $\vec{e_{loc}}$ correspond au champ appliqué $\vec{E_a}$ + champ microscopique crée par les constituants du milieu entourant $M_j$.

    $\alpha_j$ est une constante positive.
}
\subsection{Bilan}
\begin{itemize}
    \item Polarisation électronique : Modification de la répartiiton des charges internes à haque atome ou ion, toujours présente
    \item Polarisation atomique : Déplacement des atomes/ions par rapport à leur posiuton d'équilibre
    \item Polarisation d'orienrttaion : Orientation du moment dipolaire permananr d'une molcéule polaire. Nécessite une liberté de mouvement. Cela concerne donc les fluides.
\end{itemize}
\subsection{Modélisation}
\subsubsection{Polarisabilité électronique}
On utilise le modèle de Mossoti.
\begin{theorem}[Polarisabilité $\alpha_e$ selon Mossoti]
    $\alpha_e = 4\pi a^3 = \frac{Z^2z^2}{\epsilon_0m_e\omega_0^2}$ avec $\omega_0^2 = \frac{K}{m_e}$
\end{theorem}
$\alpha_e$ augmente en descendant et vers la gauche.
\subsubsection{Polarisabilié atomaiqye et ionique}
\begin{theorem}
    $$\alpha_a = \alpha_i = \frac{q^2}{\epsilon_0\mu \omega_0^2}$$
\end{theorem}
\subsubsection{Polarisabilité d'orientation $\alpha_{or}$}
Cela ne s'applique qu'aux fluides.
\begin{theorem}[$\alpha_{or}$][PASSER]
    $\alpha_{or} = \frac{p_0^2}{3\epsilon_0k_BT}$ avec $p_0$ le moment dipolaire permanent du fluide.
\end{theorem}
\section{Champ local - champ macroscopique}
On rappelle que généralement, $\vec{e_{loc}}\neq \vec{E}$
\subsection{Gaz et milieux dilués}
Dans le cas d'un gaz sous faible pression, on néglige l'intéraction entre dipole sous l'échelle microscopique, on a donc 
$\vec{e_{loc}}\simeq \vec{E}$ pour les milieux dilués.
\subsection{Champ local de Lorenz}
\begin{theorem}[PASSER]
    $\vec{e_{loc}}=\vec{E_{lorenz}}+\vec{e_{<}}$ avec $\vec{E_{lorenz}} = \vec{E}+\frac{\vec{P}}{3\epsilon_0}$

    Dans le modèle, $\vec{e_{<}}$ devient macroscopique.
\end{theorem}
\section{Lien entre polarisabilité et permitivité diélectrique}
\subsection{Gaz dilués}
\begin{theorem}[Relation de Langevin Debye]
    $$(\epsilon_r-1)\frac{M}{\rho^*}=N_A\alpha$$ avec $\rho^*$ la masse volumique et $\alpha = \alpha_e+\alpha_a+\alpha_{or}$
\end{theorem}
\subsection{Gaz denses et liquides}
\begin{theorem}[Relation de Clausius-Mossoti PASSER]
    $$\frac{\epsilon_r-1}{\epsilon_r+2}\frac{M}{\rho^*}= \frac{N_A \alpha}{3}$$ avec $\alpha = \alpha_e+\alpha_a$
\end{theorem}
\subsection{Milieux solides}
On retrouve la relation de clausisu-Mosotti.

Les formules ci-dessus ne sont pas à savoir.
\section{Aspects dynamiques}
On considère un milieu soumis à un champ électrique sinusoidal de pulsation $\omega$.
\subsection{Modèle de Drude-Lorenz}
\begin{theorem}[Modèle de Drude]
    $$m_e\frac{\dd^2 \vec{r}}{\dd t^2}=-K\vec{r}-\frac{m_e}{\tau}\frac{\dd \vec{r}}{\dd t}-e(\vec{E}+\frac{\dd \vec{r}}{\dd t}\wedge \vec{B})$$
\end{theorem}
\subsection{Susceptibilité et permittivité complexe}
On repasse en macro.
\begin{itemize}
    \item $\omega_0$ est la pulsation propre
    \item $\omega_p^2 = \frac{n_ve^2}{m_e\epsilon_0}$ la force d'oscillateur
    \item $\tau$ caractérise l'amortissement. Si $\tau >> 1$, pas d'amortissement
    \item La partie imaginaire est toujours positive
\end{itemize}
\subsection{Lien entre absorption des ondes et Partie imaginaire de er}
On se place à $\omega = \omega_0$

on utilise les valeurs réelles et non les valeurs complexes\marginInfo{En effet, on effectu un produit et le produit des parties réelles n'est pas la partie réelle du produit}
\begin{theorem}
    $$<P_m>_T = \frac{V}{2}\omega_0 \epsilon_0 (Im(e_r))E_0^2$$
\end{theorem}
On remarque que s'il il n'y a pas de frottement, il n'y a pas d'absorption.
\section{Bilan}
les différents mécanismes s'ajoutent mais ont des pulsation différentes.
\begin{itemize}
    \item Polarisabilité : toujours présente
    \item Polarisabilité ionique et atomique  : dépend de la polarité de la liaison ($\varnothing$ si liaison covalente)
    \item Polarisabiluté d'orientation : mécanisme différente, pas de résonance les collisions dominent.
\end{itemize}
\chapter{Ondes électromahnétiques dans un milieu matériel}
\section{Équations de Maxwell}
\subsection{Équations de Maxwell dans un milieu matériel}
On retrouve les équations de Maxwell. On explicite simplement $\rho = \rho_{libre}+\rho_P+\rho_{autres}$ et $\vec{j} = \vec{j_{libre}}+\vec{j_P}+\vec{j_M}$

avec 
\begin{itemize}
    \item $\rho_P = -\Div \vec{P}$
    \item $\vec{j_P} = \frac{\partial \vec{P}}{\partial t}$
    \item $\vec{j_M} = \rot \vec{M}$
\end{itemize}
En introduisant $\vec{D}$ et $\vec{H}$, on a dans un milieu matériel
\begin{itemize}
    \item $\rot \vec{H} - \frac{\partial \vec{D}}{\partial t}=\vec{j_{libre}}$
    \item $\rot \vec{E}+\frac{\partial \vec{B}}{\partial t} = \vec{0}$
    \item $\Div \vec{D} = \rho_{libre}+\rho_{autres}$
    \item $\Div \vec{B} = 0$
\end{itemize}
\subsection{Relations de passage}
On utilise un cylindre au travers de la surface et on fait tendre le flux vers 0. 
Le flux de D au travers de la surface latérale vaut donc 0.
On transforme enfin les intégrales en $\Delta S$

%TODO vérfier sur les diapositives que les relations de passage sont bien les bonnes

%Attention, on n'a pas besoin de refaire les démos, pour le cours jusqu'ici diapostive 5.6 pour le programme de révision du CC2
\section{Grandeurs énergétiques}
Dans le vide :
\begin{itemize}
    \item Densité volumique dénergie em : $w = \frac{\epsilon_0 E^2}{2}+ \frac{B^2}{2\mu_0}$
    \item Vecteur de Poyting : $\Pi = \vec{E}\wedge \frac{\vec{B}}{\mu_0}$ défini à partir du champ des réels 
    \item Pour une onde progressive dans la direction u : $\vec{Pi} = wc\vec{u}$ : Pi est un courant volumique d'énergie em.
\end{itemize}

Dans le milieu matériel linéaire :
\begin{itemize}
    \item Densité volumique dénergie em : $w = \frac{\vec{D}\cdot \vec{E}}{2}+\frac{\vec{H}\cdot \vec{B}}{2}$
    \item Vecteur de Poyting : $\Pi = \vec{E}\wedge \vec{H}$ défini à partir du champ des réels 
\end{itemize}
On introduit la puissance moyenne rayonnée à travers une surface S :
$$P = \iint_S <\vec{\Pi}>_{temps}\cdot \vec{\dd S}$$
L'intensité vaut alors $ I = \frac{P}{S}$
\section{Ondes électromagnétiqus dans un milieu Diélectrique}
\subsection{Ondes monochormatiques dans un milieu diélectrique LHI}
On consdidère des milueux localment neutres, isolants et non magnétique :
\begin{itemize}
    \item $\rho_{libre} =0$
    \item $\vec{j_{libre}} = 0$
    \item $\vec{M} = \vec{0}$
    \item $\vec{B} = \mu_0\vec{H}$
\end{itemize}
On introduit la relation constitutive du milieu : $\vec{\underline{D}} = \epsilon_0\underline{\epsilon_r}(\omega)\vec{\underline{E}}$ er $\epsilon_r = 1+ \chi_e(\omega)$
\subsubsection{$\epsilon_r$ réel et positif}
Les équations de Maxwell sont identiques à l'exception que $\epsilon_0\to \epsilon_0\epsilon_r$ :Les ondes em qui peuvent se propager dans le milieu sont semblables à celles de propageant dans le vide.

$$\lap \vec{\underline{E}}+\frac{\omega^2}{c^2}\epsilon_r(\omega)\vec{\underline{E}}=\vec{0}$$
avec la vitesse de propagation : $\frac{c}{\sqrt{\epsilon_r}}$
\subsection{Ondes planes progressives monochromatiques}
On est dans le même type de milieu que dans la section précédente et les ondes ont la même aplitude en tout point d'un même plan, i.e. des ondes planes.

Les équations de maxwell sont de la forme :
\begin{itemize}
    \item $\vec{k}\cdot \vec{E} = \vec{k}\cdot \vec{B} = 0$
    \item $\vec{k}\wedge \vec{E} = \omega \vec{B}$
    \item $\vec{k}\wedge \vec{B} = -\omega \mu_0\epsilon_0\epsilon_r(\omega)\vec{E}$
\end{itemize}
 On peut réécrire l'eq de propagation pour une onde plane :
 $$k^2 \vec{\underline{E}} = \frac{\omega^2}{c^2}\epsilon_r(\omega)\vec{\underline{E}}$$
 
 Donc si 
 \begin{itemize}
    \item $\epsilon_r$ réel et strictement positif : milieu transparent, l'onde se propage sans atténuation de l'amplitude : $k$ réel
    \item $\epsilon_r$ dépend de $\omega$ : milieu dispersif
    \item $\epsilon_r =cst$ : milieu non dispersif, c'est le cas du vide.
 \end{itemize}
On introduit la vitesse de phase de l'onde :
$v_{\varphi}(\omega) = \frac{\omega}{k} = \frac{c}{\sqrt{\epsilon_r(\omega)}}$ : La vitesse de phase est plus petite que $c$.

On peut écrire $n(\omega) = \sqrt{\epsilon_r}$.

Dans un milieu matériel, $\lambda = \frac{2\pi}{k} = \frac{\lambda_0}{n(\omega)}$

Le nombre d'onde vaut alors $ k = \frac{\omega}{c}\sqrt{\epsilon_r(\omega)}$
\subsection{Aspects énergétiques}
$$\vec{\Pi} = n(\omega)c\epsilon_0E^2\vec{u}$$
\subsection{Milieu diélectrique LHI non magnétique, non chargé, isolant, transparent et non dispersif}
\begin{itemize}
    \item $\omega = \frac{ED}{2}+\frac{B^2}{2\mu_0} = \frac{\epsilon E^2}{2}+\frac{B^2}{2\mu_0} = \epsilon_0n^2E^2$
    \item $\vec{\pi} = \frac{c}{n}w\vec{u}$
    \item %TODO finir de noter
\end{itemize}
\subsection{Milieux absorbants}
On a $\epsilon_r$ complexe donc $\vec{k}$ complexe.

Eq de maxwell, de propagation et de dispersion dans ce cas.

On note $\vec{k} = (k'+ik'')\vec{u}$

On a $\vec{E} = \vec{E_0}e^{-k''z}e^{i(k'z-\omega t)}$


L'indice optique est aussi complexe : $\underline{n} = n'+in''$

On a $\epsilon_r' = n'^2-n''^2, \epsilon_r''=2n'\cdot n''$

On a $\underline{n} = \frac{c}{\omega}\underline{k}$

On a $$<\vec{\Pi}> = \frac12 n' c\epsilon_0E_0^2e^{-2k''z}\vec{e_z}$$

Si $\epsilon_r'>0$, il y a propagation sans atténuation, milieu transparent, dans le cas contraire, il n'y a plus de propagation, que de l'atténuation.

Si $\epsilon_r = 0$

Dans ce cas, il n'y a pas de composante B, donc l'ondes est purement électrique.

Comme $\vec{k}\wedge \vec{E} = \vec{B} = 0$, l'onde n'est plus transverse mais longitudinale.

\subsection{Réflexion et transmission d'une onde à l'interface entre 2 diélectriques}
\subsection{Position du problème}
On choisit 2 milieux diélectriques LHI non dispersifs, non absorbants, non magnétiques, non chargés et très bon isolants.7
\subsection{Relations de passage}
On utilise les relations de passagepour relier les 2  champs du milieu 1 et 2.
\begin{itemize}
    \item $\vec{n_{12}}\cdot (\vec{D_2}-\vec{D_1})=0$
    \item $\vec{n_{12}}\cdot (\vec{B_2}-\vec{B_1})=0$
     \item $\vec{n_{12}}\wedge (\vec{E_2}-\vec{E_1})=0$
    \item $\vec{n_{12}}\wedge (\vec{H_2}-\vec{H_1})=0$
\end{itemize}
Il y a continuité à l'interface des composantes normales de $\vec{D},\vec{B}$ et des composantes tangentielles de $\vec{D}, \vec{E}$
\subsection{Relation de Snell-Descartes}
Les vecteurs do'nde $\vec{k_r}, \vec{k_t}$ sont contenus dans le plan d'incidence déifni par le vecteur d'onde $\vec{k_i}$ et la normale à la surface.
\subsection{Formules de Fresnel}
La polarisation est la dirrection du champ électrique

Il faut envisager 2  cas :
\begin{itemize}
    \item Polarisation contenue dans le plan d'incidence
    \item Polarisation perpendiculaire au plan d'incidence
\end{itemize}
\subsubsection{Dans le plan de l'onde}
\subsubsection{Perpendiculaire au plan d'incidence}
\subsubsection{Angle de Brewster}
\subsubsection{Coefficients de réflexion et transmissions en intensité en polarisation}
Astuce : R+T = 1
\chapter{Aspects microscopique}
\section{Moment magnétiques atomiques}
\subsection{Moment magnétique orbital}
\begin{definition}[Moment magnétique d'un électron autour d'un noyau]
    Moment d'un électron autour d'un noyau : $\vec{L} = \vec{r}\wedge (m_e \vec{v})$

    Moment d'un courant circulaire : $\vec{\mu} = iS\vec{n} = i\pi r^2\vec{n}$

    Le moment magnétique total : $\vec{\mu} = -\frac{e}{2m_e}\vec{L} = \gamma_e \vec{L}$
\end{definition}
\begin{definition}[Magnéton de Bohr]
    $\mu_{l,z} =\gamma_e L_z =  -m_l \mu_B$ avec $\mu_B$ le magnéton de Bohr.
\end{definition}
\subsection{Moment magnétique de Spin}
On ajoute le moment cinétique intrinsèque de l'électron : Spin S : $S_z = m_s \hbar$
\subsection{Moment total}
$\vec{j} = \vec{L} + \vec{S}$ et $\norm{J}^2 = (j(j+1))\hbar^2$
\section{Diamagnétisme}
On étudie le phénomène d'induction à l'échelle microscopique en présence d'un champ magnétique appliqué : $\vec{B_a} = B_a \vec{e_z}$

L'orientation de l'aimantation est opposée à celle de $\vec{B_a}$, conformément à la Loi de Lenz.

Le diamagnétisme est extrèmenent faible : $\chi_m \simeq 10^{-6}$

Dans ce cas, $\vec{B_{local}}\simeq \vec{B_a}$

Pas important
\section{Paramagnétisme}
Ici, $\chi_m \simeq 10^{-3}$


Il faut un moment cinétique total $\vec{j}\neq \vec{0}$

Le Paramagnétisme dépend de la température
\subsection{Théorie de Brillouin}
\subsubsection{Théorie}
À retenir
\subsubsection{Cas simple : 1 é célibataire}
On a $\vec{M} = n_v \mu_b \th(X)\vec{e_z}$

Quand $X>> 1, \th(X)\to 1$ : Tius les moments sont alignés sur $\vec{B_a}$. C'est le cas pour un régime en basse température ou un champ très fort.

Quand X est faible, on a $\th(X) \simeq X$ : on a une expression linéaire.

Si $n_v$ est trop grand, l'approximation des moments indépendants cesse d'être valide
\subsubsection{Gas général}
\begin{definition}
    Aimantation volumique : $\vec{M} = n_v g\mu_B B_j(X)\vec{e_z}$

\end{definition}
\section{ferromagnétisme}
Cela résulte de l'intéracion entre atomes du fait des spins des électrons des couches périphériques in complètes associées à des
sites voisins.

On peut avoir une aimantation meme en l'absence de champ appliqué. On qualifiera cette aimantation de spontannée.

\subsection{Courbe de première aimantation et cycle histérésis}
M ne prend pas les mêmes valeurs durant l'évolution de H. C'est un effet d'histérésis.

On écrit : $\mu_r = \frac{B}{\mu_0H}$
\subsection{Interpération microscopique}
\begin{definition}
    Modèle de Weiss : $\vec{B} = \vec{B_a}+\lambda \vec{M}$
\end{definition}
\end{document}
